\section{El flujo del context: el mecanismo de fondo de Elys}

Tristan insiste varias veces en que el núcleo de Elys es el flujo del context. El context no son las etiquetas de intereses del internet tradicional, sino las experiencias de una persona, su forma de expresarse, sus valores, su sentido estético, sus relaciones, sus obras pasadas y su intención actual. Elys busca que el usuario entregue ese context a su doble digital, y que el doble salga por iniciativa propia al entorno social a expresarse, comentar, recomendar y romper el hielo, para que al final traiga de vuelta conexiones entre personas reales.

\begin{figure}[H]
\centering
\includegraphics[width=0.94\textwidth]{context-flow.png}
\caption{El ciclo cerrado del flujo del context en Elys.}
\end{figure}

\begin{knowledgebox}{Lectura de la figura: el context no es una página de datos estática}
En la figura, el context circula entre el usuario, el doble cibernético, las acciones de la IA, las conexiones entre personas reales y la retroalimentación. No es una etiqueta de baja dimensión como «me gusta la música», sino un estado de alta dimensión que permite que el doble se parezca a ti, hable en tu nombre e identifique por ti oportunidades de relación.
\end{knowledgebox}

\subsection{Por qué el internet tradicional se basa en etiquetas de baja dimensión}

Las redes sociales del internet tradicional suelen depender del apodo, la foto de perfil, la región, la edad, los intereses, las relaciones de seguimiento y los registros de comportamiento. Los sistemas de recomendación pueden distribuir contenido con base en estas etiquetas, pero difícilmente entienden los valores, el sentido del humor, la forma de expresarse, la experiencia de vida y los objetivos de relación de una persona. El audio largo, los textos extensos, los historiales de chat, las obras y la retroalimentación social real apenas en la era de la IA empiezan a tener la posibilidad de entenderse de manera unificada.

\begin{importantbox}{El context es el activo de las redes sociales con IA}
La inteligencia puede volverse gradualmente accesible para todos por igual, pero el context no. Quien posea un context de usuario de más largo plazo, más auténtico y cuyo uso pueda autorizarse, estará más cerca del activo central de una red social con IA.
\end{importantbox}

\subsection{El volante de inercia del context}

El usuario aporta un context inicial y el doble empieza a actuar; las acciones generan interacciones, y las interacciones producen nueva retroalimentación; esa retroalimentación, a su vez, hace que el doble se parezca más al usuario. Lo más difícil de este volante de inercia es el arranque en frío: al principio el usuario no confía en el producto y no quiere entregar su privacidad; el producto tiene que mostrarle los beneficios reales que aporta el context antes de que lo abandone.

\begin{warningbox}{Los puntos frágiles del volante de inercia del context}
Si el usuario entrega muy poco context, el doble no se parece a él y el producto no le dice nada; si el producto pide demasiado, el usuario se preocupa por su privacidad y por ser representado de forma equivocada. La dificultad de Elys consiste justamente en encontrar un punto de arranque entre los beneficios y los riesgos.
\end{warningbox}

\subsection{Resumen de la sección}

La base de Elys no es el chat, sino el ciclo del context. Que las redes sociales con IA funcionen depende de que el context pueda obtenerse de forma segura, fluir de manera eficaz y convertirse en conexiones entre personas reales.
